\section{Periodic reduction to a finite quadratic-field maximum}
\label{sec:periodic}

Assume from now on that both $(a_i)$ and $(b_i)$ are ultimately periodic.
After increasing the preperiod, choose a common aligned period $L$.  The state
\[
 S_n=(q_n,q_{n-1},\rho_n)^T
\]
satisfies
\begin{equation}
 S_{n+1}=T(a_{n+1},b_{n+1})S_n,
 \qquad
 T(a,b)=
 \begin{pmatrix}
 a&1&0\\
 1&0&0\\
 b&0&1
 \end{pmatrix}.
 \label{eq:state-recurrence}
\end{equation}

\begin{lemma}[Three spectral scales]
\label{lem:three-scales}
Let $M$ be the product of the matrices in \eqref{eq:state-recurrence} over one
aligned tail period.  Then
\begin{equation}
 M=\begin{pmatrix}P&0\\v^T&1\end{pmatrix},
 \qquad
 \det P=(-1)^L,
 \label{eq:block-period}
\end{equation}
where $P$ is a primitive nonnegative $2\times2$ continuant matrix.  If
$\Lam>1$ is its Perron root and
$\mu=(-1)^L/\Lam$, then
\[
 \Spec(M)=\{\Lam,1,\mu\},
 \qquad |\mu|<1.
\]
For each fixed tail phase, every integer affine form in a bounded collection
of nearby $q$- and $\rho$-coordinates has an exact expansion
\begin{equation}
 C_\Lam\Lam^k+C_1+C_\mu\mu^k,
 \qquad C_\Lam,C_1,C_\mu\in\Q(\Lam).
 \label{eq:three-scale-expansion}
\end{equation}
All three coefficients are effectively computable.
\end{lemma}

\begin{proof}
The upper-left block of $T(a,b)$ is
$\left(\begin{smallmatrix}a&1\\1&0\end{smallmatrix}\right)$ and its upper-right
column is zero, so multiplication gives \eqref{eq:block-period}.  The
continuant block is primitive: every partial quotient is positive, and a
positive power of the product has all entries positive.  Each factor has
determinant $-1$, hence $\det P=(-1)^L$.

The eigenvalues of $P$ are $\Lam$ and
$\mu=\det(P)/\Lam$.  They are distinct from each other and from $1$, since
$\Lam>1$ and $|\mu|<1$.  The block form gives the stated spectrum of $M$ and
makes $M$ diagonalizable over $\Q(\Lam)$.  For example, its dominant spectral
projector is
\[
 \Pi_\Lam=
 \frac{(M-I)(M-\mu I)}{(\Lam-1)(\Lam-\mu)}.
\]
The other two projectors are obtained cyclically.  If $A_j$ is the fixed
prefix product leading to phase $j$, then
$S_{n_0+j+kL}=A_jM^kS_{n_0}$.  Applying an integer linear functional and the
three projectors gives \eqref{eq:three-scale-expansion}; an affine constant is
absorbed into $C_1$.  Nearby indices only change the fixed phase product.
\end{proof}

\begin{remark}
\label{rem:not-primitive}
The full matrix $M$ is reducible: its third column is fixed.  Perron--Frobenius
theory is used only for the continuant block $P$.  No primitivity assertion
about the $3\times3$ state matrix is needed or valid.
\end{remark}

The next lemma contains the finite-mass step of the proof.

\begin{lemma}[At most four candidates per phase]
\label{lem:four-per-phase}
Fix a tail phase.  Every nonempty row of \cref{tab:eight-cases} has a
repetition denominator whose coefficient of $\Lam^k$ is positive.  The
right-endpoint ratio therefore converges to an effectively constructible
element of $\Q(\Lam)$.  There are at most four such limits in the phase.
\end{lemma}

\begin{proof}
Along a fixed phase, both $q_n$ and $q_{n-1}$ grow as positive multiples of
$\Lam^k$.  It remains to exclude cancellation in the other displayed
denominators.  This is transparent row by row from
$0\leq s<p$, $0\leq t<q$, and $u=t+bq$.

In cases~1--3 the only additional forms are $Q-s$ and $Q-t$; in case~3,
$a\geq2$.  They dominate $q$ or $p$.  In case~4,
$Q-t=q+p-t>p$.  In case~5, $b\leq a-2$, and hence
\[
 Q-u=(a-b)q+p-t>q,
 \quad Q-bq=(a-b)q+p\geq2q+p,
 \quad Q+q-u>2q.
\]
In case~6, $t=s<p$, and its forms $q+p-s$, $q+p$, and $2q+p-s$
dominate $q$.  In case~7, $q+p-t>p$ and
$2q+p-t>q+p$.  Finally, case~8 uses $p$ and $q+p-s>q$.
Thus every displayed repetition value is bounded below by either $p$ or $q$
on its phase and has positive dominant coefficient.

By \cref{lem:three-scales}, both the right endpoint $B_k$ and the denominator
$V_k$ have expansions \eqref{eq:three-scale-expansion}.  Since the dominant
coefficient of $V_k$ is positive,
\[
 \lim_{k\to\infty}\frac{B_k}{V_k}
 =\frac{C_\Lam(B)}{C_\Lam(V)}\in\Q(\Lam).
\]
The finite constants $-1$ and $-2$ disappear only at this step.  Inspection of
\cref{tab:eight-cases} gives at most four rows in any one phase.
\end{proof}

For completeness, the phase maximum can already be stated independently of
the algorithm used to decide empty rows.  If $\mathcal R_j$ is the set of rows
that are eventually nonempty in phase $j$, then
\begin{equation}
 \dio(x)=1+
 \max_{0\leq j<L}\ 
 \max_{R\in\mathcal R_j}
 \frac{C_\Lam(B_{j,R})}{C_\Lam(V_{j,R})}.
 \label{eq:finite-maximum}
\end{equation}
Indeed, \cref{prop:finite-interface} reduces each plateau to its right
endpoint, every phase sequence converges by \cref{lem:four-per-phase}, and the
limsup of finitely many convergent phase sequences is the maximum of their
limits.  The remaining issue---deciding $\mathcal R_j$ and producing certified
algebraic output---is exact and finite.

